% diversity_reduction_synthesis.tex [NEW]
% Connects the local (Taylor) and global (geometric) analyses,
% explains the scaling correspondence, and summarizes the full
% theoretical picture.
%
% Intended to follow all three previous files as a concluding subsection.
%
% === NOTES ===
%
% This file ties the three proof components together:
%   1. Local/Taylor (diversity_reduction_theory.tex): covariance structure
%   2. Global/geometric (diversity_reduction_geometric.tex): exact bounds
%   3. Conditions (diversity_reduction_conditions.tex): when it holds
%
% The scaling correspondence (O(1/||s||) chordal vs O(1/||s||^2) variance)
% is not a contradiction -- it follows from the standard relationship
% between angular distance and variance for concentrated distributions.
%
% === END NOTES ===

\subsection{Synthesis}
\label{sec:synthesis}

We now summarize the full theoretical picture and clarify the relationship
between the local and global analyses.

\subsubsection{Scaling Correspondence}

The global analysis (Section~\ref{sec:geometric-bound}) gives bounds
that shrink as $O(1/\|s\|)$: spherical variance, expected pairwise
chordal distance, and diameter all scale inversely with the steering
magnitude.
The local analysis (Section~\ref{sec:covariance}) gives total variance
$\mathrm{tr}(\Sigma_z) = O(1/\|s\|^2)$.
These are consistent, not contradictory.

For a distribution concentrated near a point on the sphere, chordal
distance $c$ and angular distance $\theta$ are related by
$c = 2\sin(\theta/2) \approx \theta$ for small $\theta$.
The variance of a scalar angular displacement scales as $\theta^2$,
while the angular displacement itself scales as $\theta$.
So $O(1/\|s\|)$ in chordal distance corresponds to $O(1/\|s\|^2)$
in variance, which is exactly the relationship between our two bounds.

More precisely, for a distribution on $S^{d-1}$ concentrated near a
pole with small angular spread $\theta$, the spherical variance
$V = 1 - \|\bar{R}\| \approx \theta^2/2$ and the trace of the
tangent-space covariance $\mathrm{tr}(\Sigma_z) \approx \theta^2$.
The global bound gives $\theta = O(1/\|s\|)$; squaring recovers the
local bound.

\subsubsection{What Each Result Provides}

The two analyses are complementary, and each answers questions the other
cannot:

\paragraph{The local (Taylor) analysis} reveals the internal structure
of the diversity reduction.
It identifies \emph{which} dimension of variation is lost (the direction
$\hat{m} = \widehat{\mu + s}$), gives the full eigenvalue structure of
the output covariance, and shows that effective rank drops by exactly one
when $\Sigma_x$ has nonzero variance along $\hat{m}$.
This structural detail connects directly to the pass@k story: the lost
dimension corresponds to a specific mode of computational variation that
the model can no longer express.
The cost is that the result is approximate, valid only when input
variation is small relative to $\|\mu + s\|$.

\paragraph{The global (geometric) analysis} provides exact,
assumption-free bounds on the absolute spread of post-normalization
representations.
It requires no distributional assumptions beyond bounded support, and
the bounds hold for any steering magnitude.
The cost is that it provides only scalar bounds (diameter, spherical
variance), not structural information about which dimensions are
compressed.

\paragraph{The condition analysis} (Section~\ref{sec:conditions})
identifies when the diversity reduction is genuine relative to the
unsteered baseline: $\|\mu + s\| > \|\mu\|$.
This is the load-bearing assumption that connects the absolute
concentration bounds to relative diversity claims.

\subsubsection{The Complete Argument}

Combining all three components, the complete theoretical argument is:

\begin{enumerate}
  \item \textbf{Architecture.}
    Every transformer block contains normalization layers (RMSNorm or
    LayerNorm) that project the residual stream onto a sphere (or
    affine image thereof).

  \item \textbf{Mechanism.}
    Adding a fixed steering vector $s$ shifts the pre-normalization
    distribution.
    The normalization map then concentrates all representations toward
    the pole $\hat{s}$, with chordal spread bounded by
    $O(\mathbb{E}[\|x\|] / \|s\|)$
    (Proposition~\ref{prop:concentration}).

  \item \textbf{Structure.}
    The concentration specifically eliminates variation along the
    direction $\hat{m} = \widehat{\mu + s}$ and compresses all
    remaining variation by $1/\|\mu + s\|^2$
    (Proposition~\ref{prop:diversity-reduction}).

  \item \textbf{Condition.}
    This constitutes a \emph{reduction} relative to the unsteered
    baseline when $\|\mu + s\| > \|\mu\|$
    (Section~\ref{sec:conditions}), a condition satisfied generically
    in practice.

  \item \textbf{Compounding.}
    The steering vector persists via residual connections, so the
    compression applies at every subsequent normalization layer.
    The net effect per layer is normalization-induced compression
    minus sub-layer-induced expansion; empirical evidence
    \citep{sinii2025small} confirms that the normalization step is the
    specific bottleneck limiting steering vector effectiveness.

  \item \textbf{Output.}
    Representational diversity upper-bounds output diversity, since the
    unembedding is a linear map that cannot increase rank.
\end{enumerate}

\noindent
Steps 1--3 are proved.
Step 4 is a stated condition with strong practical justification.
Step 5 is a qualitative argument supported by empirical evidence.
Step 6 is exact.

\subsubsection{Empirical Predictions}

The theory makes several testable predictions:

\begin{enumerate}
  \item \textbf{Effective rank decreases with steering magnitude.}
    Measuring the effective rank of the representation covariance at
    each layer, steered models should show lower effective rank than
    unsteered models, with the reduction scaling as
    $O(1/\|s\|^2)$ in variance terms.

  \item \textbf{The lost dimension aligns with the steering direction.}
    The principal component of variance lost under steering should
    have high cosine similarity with $\hat{m} = \widehat{\mu + s}$.

  \item \textbf{Diversity reduction concentrates at normalization layers.}
    Measuring representational diversity before and after each
    sub-component within a transformer block (following the
    methodology of \citet{sinii2025small}), the largest drops should
    occur at the normalization steps, not at the attention or MLP
    steps.

  \item \textbf{Pass@k degrades faster than pass@1.}
    If steering compresses the output distribution while preserving
    the mode, then pass@1 (which depends only on the mode) should
    be relatively preserved while pass@k for large $k$ (which depends
    on distributional coverage) should degrade.
    This prediction connects the representation-level theory to the
    behavioral-level measurements in the experimental sections.
\end{enumerate}
